\section{Aufbereitete Engineering-Testdaten der Verifikation und Validierung}
\label{app:vv_test_data}

Die für die formative End-to-End- und Multi-Source-Verifikation verwendeten
Testdaten wurden aus realen Engineering-Informationen von PreciPoint
abgeleitet. Inhalt und fachliche Zusammenhänge orientieren sich an einem
realen Remote-Collaboration-Anwendungsfall aus der digitalen Mikroskopie.

Für die kontrollierte Evaluation wurden die zugrunde liegenden Informationen
hinsichtlich Umfang, Struktur und Formulierung gezielt reduziert und
aufbereitet. Ziel war es, die relevanten fachlichen Zusammenhänge und
Informationsüberschneidungen zwischen mehreren Dokumenttypen zu erhalten und
gleichzeitig einzelne Verarbeitungsschritte eindeutig nachvollziehbar zu
machen. Die Begrenzung des Informationsumfangs reduzierte zudem Laufzeit und
Tokenbedarf der wiederholten LLM-basierten Testläufe.

Die nachfolgend dokumentierten Dateien stellen daher keine unveränderten
Produktivdokumente des Unternehmens dar. Ihre fachliche Grundlage stammt
jedoch aus realen Engineering-Daten. Die vier Dokumente bilden unterschiedliche
Informationsperspektiven auf denselben Produktkontext ab: Produktbeschreibung,
User Workflow, System Requirements und technische Architektur.

\subsection{Product Overview}

\lstinputlisting[
    basicstyle=\ttfamily\footnotesize,
    breaklines=true,
    columns=fullflexible
]{Anhang/Testdaten_WP12/01_product_overview.md}

\subsection{User Workflow Notes}

\lstinputlisting[
    basicstyle=\ttfamily\footnotesize,
    breaklines=true,
    columns=fullflexible
]{Anhang/Testdaten_WP12/02_user_workflow.md}

\subsection{Draft System Requirements}

\lstinputlisting[
    basicstyle=\ttfamily\footnotesize,
    breaklines=true,
    columns=fullflexible
]{Anhang/Testdaten_WP12/03_system_requirements.md}

\subsection{Technical Architecture Notes}

\lstinputlisting[
    basicstyle=\ttfamily\footnotesize,
    breaklines=true,
    columns=fullflexible
]{Anhang/Testdaten_WP12/04_technical_architecture_notes.md}